\begin{proposition}[G-16: Explicit nonrectangular labels throughout the dimension-certified range]\label{G-16}
Let $n>m\geq3$, put $a=n-m$, and suppose
\begin{equation}\label{G16:dimensionhypothesis}
 \dim\mathcal P_{m,4}+3>2n^2+2.
\end{equation}
Define
\[
 b=\binom{n+3}{3}\binom{a+3}{3},\qquad L=n+3a,
 \qquad \beta=6bL^{30},\qquad u=2n^2-10,
\]
and set
\begin{equation}\label{G16:firstdegree}
 C=((n+a)\beta^2)^u,\qquad
 s_*=(u+1)(C-1)+1,\qquad d_* =\beta s_*.
\end{equation}
There is an integer $1\leq d\leq d_*$ with $4\mid md$ such that,
for the explicit shape
\begin{equation}\label{G16:label}
 \lambda_{m,n,d}=
 \left(ad+\frac{md}{4},\frac{md}{4},\frac{md}{4},\frac{md}{4}\right),
\end{equation}
one has
\[
 m_{\lambda_{m,n,d}}(\ell^a\operatorname{per}_m)
       >m_{\lambda_{m,n,d}}(\det_n).
\]
The bound has $\log d_*=O(n^2\log n)$ uniformly for $3\leq m<n$.
There is no additional condition $n<4m/3$ in this proposition.

An explicit progression of successful degrees is also available. Put
\[
 T=\binom{b+\beta}{b}-1,\quad E=4T(\beta!),\quad
 Q=((n+a)E)^{u-1},\quad K=u(Q-1)+1.
\]
For every integer $k\geq K$, the same inequality holds at
$d=Ek$ and the label \eqref{G16:label}.
Because $a>0$, every displayed label is nonrectangular.
\end{proposition}

\paragraph{Scope and sources.}
This applies to every padded pair certified by R-14, R-16, R-18,
R-19, or R-20, including those at or above $n=4m/3$, where G-12
excludes rectangular gaps. The construction uses a Segre--Veronese
coordinate ring and the same Kempf--Ness and reductive-quotient
facts used in G-11. For the quantitative bounds, we use Derksen
(2001), Theorem~1.1 and Proposition~1.2, printed p.~956, and the
image-degree/Chardin argument of G-10 and G-13. The new representation
requires a different coefficient-degree bound, which is checked below.

\begin{proof}
Let $V=\mathbb C^4$, $G=\mathrm{SL}(V)$, and put
\[
 Y_D=\mathcal D_{n,4},\qquad
 Y_P=\{L^a q:L\in V,\ q\in\mathcal P_{m,4}\}^{\rm cl}.
\]
These are the irreducible cones used in R-14, with
\[
 \dim Y_D=2n^2+2,\qquad
 \dim Y_P=\dim\mathcal P_{m,4}+3.
\]
For either cone $Y\subseteq\operatorname{Sym}^n V$, define
\begin{equation}\label{G16:enhancedcone}
 X_Y=\{F\otimes v^a:F\in Y,\ v\in V^*\}^{\rm cl}
 \subseteq W:=\operatorname{Sym}^n V\otimes\operatorname{Sym}^a V^*.
\end{equation}

\paragraph{1. The graded ring and the selected representation.}
The projectivization of \eqref{G16:enhancedcone} is the closed
Segre--Veronese image of $\mathbb P(Y)\times\mathbb P(V^*)$.
In particular $X_Y$ is irreducible and
\begin{equation}\label{G16:xdimension}
 \dim X_Y=\dim Y+3.
\end{equation}
Its coordinate ring has the exact grading
\begin{equation}\label{G16:segrering}
 \mathbb C[X_Y]_d=
       \mathbb C[Y]_d\otimes\operatorname{Sym}^{ad}V.
\end{equation}
Indeed each linear coordinate restricts to a product of a degree-one
coefficient function of $F$ and a degree-$a$ polynomial function of $v$.
Their degree-$d$ products span the right side: the first factor is
generated in degree one, and every degree-$ad$ monomial in $v$ can
be split into $d$ degree-$a$ monomials. The only additional relations
are exactly the restrictions defining this image. This proves
\eqref{G16:segrering} directly, without a projective-normality assumption
on $Y$.

To keep the duals explicit, use the actual coordinate-ring decomposition
\[
 \mathbb C[Y]_d=
       \bigoplus_{\lambda\vdash nd}m_\lambda(Y)S_\lambda V^*.
\]
It is the dual of the positive-weight orbit-span convention used
elsewhere in the paper. Schur's lemma for $G$ and
\eqref{G16:segrering} give
\begin{equation}\label{G16:extraction}
 \dim\mathbb C[X_Y]_d^G=
 \begin{cases}
 m_{\lambda_{m,n,d}}(Y),&4\mid md,\\
 0,&4\nmid md.
 \end{cases}
\end{equation}
For full detail,
\[
 (S_\lambda V^*\otimes\operatorname{Sym}^{ad}V)^G
   =\operatorname{Hom}_G(S_\lambda V,\operatorname{Sym}^{ad}V).
\]
This has dimension one precisely when
$\lambda=(ad+r,r,r,r)$ for an integer $r\geq0$, and otherwise is
zero. The fixed size $|\lambda|=nd$ forces $4r=md$, proving
the formula and ruling out any additional labels in this extraction.

\paragraph{2. Stable points with arbitrary padding.}
Let $x_0,x_1,x_2,x_3$ be an orthonormal basis of $V$ and
$e_0^*$ its first dual vector. Set
\[
 \gamma=\sqrt{\binom{n}{m}},\qquad
 F=x_0^a\bigl(x_0^m+\gamma x_1^m+
                         \gamma x_2^m+\gamma x_3^m\bigr),
 \qquad w=F\otimes(e_0^*)^a.
\]
The degree-$m$ factor of $F$ belongs to both
$\mathcal P_{m,4}$ and $\mathcal D_{m,4}$. In fact each of its two
binary summands factors into $m$ linear factors over $\mathbb C$;
put those factors on the diagonal and the cycle of an $m\times m$
pencil. Its permanent is the sum of the two products, and its
determinant is the same sum after absorbing $(-1)^{m-1}$ into one
cycle factor. Block padding then gives $F\in Y_P\cap Y_D$.
Thus the same tensor $w$ lies in both enhanced cones.

Use the Hermitian structures induced from tensor powers. The four
monomial tensors in $w$ are mutually orthogonal and have equal norm:
$\|x_0^n\|=1$, while
$\|x_0^a x_i^m\|^2=a!m!/n!$ for $i>0$, which is canceled by
$\gamma^2$. Their torus weights are exactly
\[
 ne_0-ae_0=me_0,\qquad
 ae_0+me_i-ae_0=me_i\quad(1\leq i\leq3).
\]
Their mean weight is $(m/4)(1,1,1,1)$, so the traceless diagonal
moment map is zero. The off-diagonal terms vanish as well. On the
first tensor factor, the four monomial weights differ by
$m(e_i-e_j)$, never by one root when $m\geq3$. On the second
factor, an off-diagonal derivative is either zero or orthogonal
to $(e_0^*)^a$. Therefore $w$ has zero $G$ moment map.

Its projective stabilizer has zero Lie algebra. To verify this
without a generic-stabilizer assertion, suppose
$A\in\mathfrak{sl}_4$ and $A\cdot F=\kappa F$.
Write $F=x_0^a q$. Dividing by $x_0^{a-1}$ and reducing modulo
$x_0$ shows that $A\cdot x_0$ is a multiple $\eta x_0$, since
$a>0$ and $q|_{x_0=0}\ne0$. It follows that
$A\cdot q=(\kappa-a\eta)q$.
For the Fermat sum $q$, the monomials $x_j^{m-1}x_i$, $i\ne j$,
are all distinct because $m\geq3$. Their coefficients force every
off-diagonal entry of $A$ to vanish. The four pure-power coefficients
then force every diagonal entry to be the same, and trace zero
forces $A=0$. If $A$ stabilizes the projective tensor line of $w$,
the elementary tangent property of a nonzero decomposable tensor
forces it to stabilize the line of $F$; hence again $A=0$.
Thus the projective stabilizer is finite.

Kempf--Ness now makes $w$ a stable point of each enhanced cone.
Their stable loci are nonempty open sets. Using
\eqref{G16:xdimension}, the two affine invariant quotients therefore
have dimensions
\begin{equation}\label{G16:quotdims}
 \begin{split}
 \dim A_D&=u=2n^2-10,\\
 \dim A_P&=\dim\mathcal P_{m,4}-9>u,
 \end{split}
 \qquad
 A_D=\mathbb C[X_{Y_D}]^G,\quad A_P=\mathbb C[X_{Y_P}]^G.
\end{equation}
The subtraction is $15=\dim G$. The padding weights were canceled
in the tensor construction, so this calculation needs no inequality
between $n$ and $4m/3$.

\paragraph{3. A quantitative invariant bound for this representation.}
The dimension of $W$ is $b$. In the matrix coordinates on
$\mathrm{SL}_4$, the action on $V^*$ is given by the transpose
adjugate, whose entries have degree three. The action on $W$
therefore has matrix coefficients of degree at most
$L=n+3a$. Its kernel is finite: its differential is a nonzero
representation of the simple Lie algebra $\mathfrak{sl}_4$, hence
injective, and a zero-dimensional algebraic kernel is finite.
For example the vector $x_0^n\otimes(e_1^*)^a$ has a nonzero
diagonal Lie-algebra weight, verifying nontriviality.

Derksen's Proposition~1.2 applies to the embedding
$\mathrm{SL}_4\subseteq\mathbb A^{16}$ defined by the degree-four
equation $\det-1=0$. Thus its parameters are $t=16$, group dimension
$15$, equation-degree bound $4$, and action-degree bound $L$, giving
$\sigma(W)\leq4L^{15}$.
Theorem~1.1, with invariant-quotient dimension at most $b$, then gives
generators of $\mathbb C[W]^G$ of degrees at most
\[
 \max\{2,(3/8)b(4L^{15})^2\}\leq6bL^{30}=\beta.
\]
Exactness of invariants for this reductive group makes $A_D$ and
$A_P$ graded quotients of that invariant ring. Choose common
homogeneous generators $g_i$ of degrees $1\leq e_i\leq\beta$.

The determinant enhanced cone has a dense parametrization
\[
 (Z,v)\longmapsto\det(Z_1x_0+Z_2x_1+Z_3x_2+Z_4x_3)\otimes v^a.
\]
Its coordinates are homogeneous of degree $n+a$ in the
$4n^2+4$ parameters. This parameter degree is distinct from the
group-action coefficient degree $n+3a$ above.
The generator coordinates $g_i$ composed with this map have degree
$(n+a)e_i\leq(n+a)\beta$.
Homogenizing them with one additional source coordinate gives a
dominant rational map to the ordinary projective closure of
$\operatorname{Spec}A_D$ in its generator embedding. The image-degree
argument of G-10, including its treatment of a base locus, yields
\[
 \deg\overline{\operatorname{Spec}A_D}
      \leq((n+a)\beta)^u.
\]

Let $H_D^{\leq}(d)=\sum_{j=0}^d\dim(A_D)_j$, and define
$H_P^{\leq}$ similarly. The graded pieces through degree $d$
are contained in the ordinary degree-at-most-$d$ generator filtration,
because all generator weights are at least one. Chardin's bound on
the ordinary projective closure therefore gives
\begin{equation}\label{G16:cumulativeupper}
 H_D^{\leq}(d)\leq((n+a)\beta)^u\binom{d+u}{u}.
\end{equation}
By \eqref{G16:quotdims}, at least $u+1$ images of the generators in
$A_P$ are algebraically independent. Their monomials of ordinary
total degree at most $s$ have graded degree at most $\beta s$;
hence
\begin{equation}\label{G16:cumulativelower}
 H_P^{\leq}(\beta s)\geq\binom{s+u+1}{u+1}.
\end{equation}
As in G-13, the exact binomial comparison is
\[
 \frac{\binom{s+u+1}{u+1}}{\binom{\beta s+u}{u}}
 =\frac{s+u+1}{u+1}\prod_{i=1}^u\frac{s+i}{\beta s+i}
 \geq\frac{s+u+1}{(u+1)\beta^u}.
\]
At $s=s_*$ this is strictly greater than
$C/\beta^u=((n+a)\beta)^u$. Equations
\eqref{G16:cumulativeupper}--\eqref{G16:cumulativelower} give a
strict cumulative gap through $d_*$. Degree zero contributes one
on both sides, so there is a positive degree $d\leq d_*$ with
$\dim(A_P)_d>\dim(A_D)_d$.
Formula \eqref{G16:extraction} identifies this gap with exactly the
label \eqref{G16:label} and forces $4\mid md$.
The four-variable highest-weight restriction of G-05/R-14 transfers
it to the full padded-permanent and determinant orbit closures.
Since $a<n$, $b=O(n^6)$, $L=O(n)$, and $\beta=O(n^{36})$;
\eqref{G16:firstdegree} consequently has
$\log d_*=O(n^2\log n)$.

\paragraph{4. An explicit progression, if individual degrees are desired.}
There are at most $T=\binom{b+\beta}{b}-1$ generators, since the
ambient polynomial space in positive degrees through $\beta$ has
that dimension. The elementary generator-grouping proof in G-11
therefore makes both $E$th Veronese rings standard graded for
$E=4T(\beta!)$. Their dimensions are still $u$ and at least $u+1$.
The degree-one invariant functions of the first Veronese, composed
with the enhanced determinant parametrization, have degree
$(n+a)E$. The projective image has dimension $u-1$ and degree at
most $Q=((n+a)E)^{u-1}$.
Chardin's theorem and linear Noether normalization give respectively
\[
 \dim(A_D)_{Ek}\leq Q\binom{k+u-1}{u-1},\qquad
 \dim(A_P)_{Ek}\geq\binom{k+u}{u}.
\]
Their binomial ratio is $(k+u)/u$, strictly larger than $Q$ for
$k\geq u(Q-1)+1$. This proves the claimed progression; $4\mid E$
ensures all labels in it are integral.
\end{proof}

\paragraph{Examples where ordinary rectangles are excluded.}
The pairs $(m,n)=(15,20)$ and $(16,22)$ from R-17 satisfy
\eqref{G16:dimensionhypothesis}. G-16 produces labels
$\lambda=(35r,15r,15r,15r)$ at degrees $d=4r$ for the first pair,
and labels $\lambda=(10d,4d,4d,4d)$ for the second, at the
successful degrees specified above. G-12 precludes ordinary
rectangular gaps for both pairs. The same explicit label family
applies to the unbounded nonrectangular strip supplied by R-20.

